% =====================================================================================================
% Kapitel 1: Einleitung                                                         Zielumfang: ca. 3 Seiten
% -----------------------------------------------------------------------------------------------------
% Legende der Kommentarblöcke (gilt für alle Kapitel):
%   ZIEL     – Was soll der Leser nach dem Abschnitt wissen?
%   INHALT   – Stichpunkte, die vorkommen sollten (Reihenfolge ist ein Vorschlag)
%   QUELLEN  – Woher die Belege kommen (Schlüssel aus literature.bib, sobald angelegt)
%   UMFANG   – grobe Seitenzahl
%   HINWEIS  – Stolperfallen
% Den Fließtext schreibt ihr selbst unter den jeweiligen Block. Die Blöcke können danach stehen
% bleiben (sie erscheinen nicht im PDF) oder gelöscht werden.
% =====================================================================================================

\section{Einleitung}

\subsection{Motivation}
% ZIEL:    Warum lohnt sich die Überarbeitung von Sicherheitskreis, Scanner-Anbindung und IO-Link-Master?
% INHALT:  - Modellanlage dient der Lehre → muss sicher, wartbar und nachvollziehbar dokumentiert sein
%          - Anlage wurde über mehrere Arbeiten erweitert → provisorische Lösungen (Relais lose am
%            Prototyp, Scanner am Laptop, IO-Link-Master auf der Schleppkette)
%          - kurzer Bezug zur Praxis: funktionale Sicherheit, vernetzte Feldgeräte, Wartbarkeit
% QUELLEN: Vorgängerarbeiten T3100/T3200 (Ist-Zustand); ggf. ein Fachbuch zur funktionalen Sicherheit
% UMFANG:  ca. 1 Seite
% HINWEIS: Kein allgemeines „Industrie 4.0“-Intro – direkt zur Anlage kommen.

\subsection{Ausgangslage und Vorarbeiten}
% ZIEL:    Leser weiß, was die Anlage heute kann und auf welchen Arbeiten diese Arbeit aufbaut.
% INHALT:  - DHBW-Modell-Hochregallager: S7-1500, TIA Portal V18, HMI, RFID, IO-Link, QR-Scanner
%          - Vorarbeit T3100 Schaal & Zorko: Lichtgitter, Umlagerung
%          - Vorarbeit T3200 Schaal & Zorko: Auftragsmanagement mit Python/Snap7 und QR-Scanner
%          - Hinweis auf Kapitel 2.1 für die Details
% QUELLEN: Vorgängerarbeiten (zitieren, eigenständig zusammenfassen!)
% UMFANG:  ca. 1 Seite

\subsection{Zielsetzung und Aufbau der Arbeit}
% ZIEL:    Klare, prüfbare Ziele + Leitfaden durch die Kapitel.
% INHALT:  - Ziel 1: Sicherheitsrelais PNOZ s4 vom Prototyp auf Hutschiene am Steuerungsgestell,
%            fachgerecht verdrahtet und geprüft
%          - Ziel 2: QR-Scanner über RS232-Ethernet-Konverter ins Netzwerk, Python auf TCP/IP
%          - Ziel 3: IO-Link-Master von der Schleppkette an eine feste Position
%          - Aufbau: je ein Satz zu Kapitel 2 bis 7
% UMFANG:  ca. 1 Seite
% HINWEIS: Ziele so formulieren, dass im Fazit (Kapitel 7) „erreicht / teilweise / nicht“ prüfbar ist.
